\chapter{Fourier Series}

\begin{center}\itshape
Profound study of nature is the most fertile source of mathematical discoveries.\\
\upshape --- Joseph Fourier, \emph{Th\'eorie analytique de la chaleur}, 1822
\end{center}

In 1807 Fourier claimed that every function on an interval, even one given by different formulas on different pieces,
is a sum of sines and cosines. Lagrange objected, and he had a point: Fourier could not say in which sense the sum
converges, and the question kept analysts busy for a century. It forced the modern definitions of function, of
integral and of uniform convergence, and it led Cantor to set theory.

This chapter answers the question in the three senses that matter for computing. Pointwise, the series of a piecewise
$C^1$ function converges to the function, or to the average of its two limits at a jump (Dirichlet). Uniformly, the
averages of the partial sums of a continuous function converge to it (Fej\'er), even though the partial sums themselves
may diverge. In mean square, the series of any piecewise continuous function converges to it, and the energy of the
function is the sum of the energies of its harmonics (Parseval). Along the way we meet the Gibbs overshoot, two more
proofs of $\zeta(2)=\frac{\pi^2}6$, Weyl's theorem on the equidistribution of $n\alpha$ modulo $1$, Fourier inversion on $\Z^d$
promised in Tome~3, and Hurwitz's proof that the circle encloses the largest area for its length.

\begin{think}
The functions $\sin x,\sin2x,\sin3x,\dots$ are all smooth. Can a sum of them have a jump? And if it has one, what value do
you expect the series to take at the jump?
\end{think}

\section{Fourier coefficients and best approximation}

Throughout, functions are $2\pi$-periodic from $\R$ to $\C$. Such a function is \emph{piecewise continuous} if $[-\pi,\pi]$ splits
into finitely many intervals on whose interiors $f$ is continuous with finite limits at both ends; its integral is the
sum of the integrals over the pieces. It is \emph{piecewise $C^1$} if moreover on each piece $f$ coincides with a $C^1$ function
on the closed interval. Since $f$ is periodic, its integral over any interval of length $2\pi$ is the same (cut the interval
at a multiple of $\pi$ and translate one part by $2\pi$).

For such functions let
\[
\begin{gathered}\langle f,g\rangle=\frac1{2\pi}\int_{-\pi}^\pi f(t)\overline{g(t)}\dd t,\qq\|f\|_2=\sqrt{\langle f,f\rangle},\end{gathered}
\]
and let $e_n(t)=e^{int}$ for $n\in\Z$. Since $\int_{-\pi}^\pi e^{ikt}\dd t=0$ for every integer $k\ne0$, the $e_n$ are \emph{orthonormal}:
$\langle e_m,e_n\rangle$ is $1$ if $m=n$ and $0$ otherwise. The form $\langle\cdot,\cdot\rangle$ is linear in its first argument, Hermitian, and
$\langle f,f\rangle\ge0$, so Cauchy--Schwarz holds with the proof of Tome~0, Chapter~4: $|\langle f,g\rangle|\le\|f\|_2\|g\|_2$. (A piecewise continuous
$f$ can have $\|f\|_2=0$ without vanishing, by changing finitely many values; this does no harm below.)

\begin{definition}
The \emph{Fourier coefficients} of $f$ are
\[
c_n(f)=\langle f,e_n\rangle=\frac1{2\pi}\int_{-\pi}^\pi f(t)e^{-int}\dd t,
\]
and its \emph{partial Fourier sums}
are $S_Nf=\sum_{|n|\le N}c_n(f)e_n$. A \emph{trigonometric polynomial of degree at most $N$} is a combination of $e_{-N},\dots,e_N$.
\end{definition}

For real $f$, grouping $n$ with $-n$ gives the real form $S_Nf(x)=\frac{a_0}2+\sum_{n=1}^N(a_n\cos nx+b_n\sin nx)$, with
$a_n=\frac1\pi\int f(t)\cos nt\dd t$ and $b_n=\frac1\pi\int f(t)\sin nt\dd t$.

\begin{proposition}[best approximation]\label{prop:bestapprox}
Let $f$ be piecewise continuous and $P$ a trigonometric polynomial of degree at most $N$. Then
\[
\|f-P\|_2^2=\|f-S_Nf\|_2^2+\|S_Nf-P\|_2^2.
\]
In particular $S_Nf$ is the trigonometric polynomial of degree at most $N$ closest to $f$, and
$\|f-S_Nf\|_2^2=\|f\|_2^2-\sum_{|n|\le N}|c_n(f)|^2$.
\end{proposition}
\begin{proof}
For $|n|\le N$, $\langle f-S_Nf,e_n\rangle=c_n(f)-c_n(f)=0$ by orthonormality, so $u=f-S_Nf$ is orthogonal to $v=S_Nf-P$, which is a
combination of these $e_n$. Then $\|u+v\|_2^2=\|u\|_2^2+\|v\|_2^2+2\operatorname{Re}\langle u,v\rangle=\|u\|_2^2+\|v\|_2^2$. With $P=0$ this reads
$\|f\|_2^2=\|f-S_Nf\|_2^2+\|S_Nf\|_2^2$, and $\|S_Nf\|_2^2=\sum_{|n|\le N}|c_n(f)|^2$ by orthonormality.
\end{proof}

\begin{corollary}[Bessel's inequality and Riemann--Lebesgue]\label{cor:bessel}
For $f$ piecewise continuous, $\sum_{n\in\Z}|c_n(f)|^2\le\|f\|_2^2$. Hence $c_n(f)\to0$ as $|n|\to\infty$, and also
\[
\int_{-\pi}^\pi f(t)\sin\Bigl(N+\frac12\Bigr)t\dd t\xrightarrow[N\to\infty]{}0.
\]
\end{corollary}
\begin{proof}
The partial sums of $\sum|c_n|^2$ are at most $\|f\|_2^2$ by Proposition~\ref{prop:bestapprox}; the terms of a convergent series tend
to $0$. For the last claim, $\sin(N+\frac12)t=\sin Nt\cos\frac t2+\cos Nt\sin\frac t2$, and $\sin Nt$, $\cos Nt$ are combinations of $e_{\pm N}$,
so the integral is a combination of $c_{\pm N}(f\cos\frac t2)$ and $c_{\pm N}(f\sin\frac t2)$ with bounded coefficients.
\end{proof}

This completes the Riemann--Lebesgue lemma of Tome~0, Exercise~6.6, which assumed $f$ of class $C^1$: here $f$ is only
piecewise continuous, and the proof is geometry (Pythagoras in a space of functions) rather than an integration by parts.

\section{Dirichlet's theorem}

\begin{lemma}[Dirichlet kernel]\label{lem:dirichletkernel}
For every $N\ge0$ and every $x$, $S_Nf(x)=\frac1{2\pi}\int_{-\pi}^\pi f(x+u)D_N(u)\dd u$, where
\[
D_N(u)=\sum_{|n|\le N}e^{inu}=\frac{\sin(N+\frac12)u}{\sin\frac u2}
\]
for $u\notin2\pi\Z$, and $D_N(0)=2N+1$. Moreover $D_N$ is even and $\frac1{2\pi}\int_0^\pi D_N(u)\dd u=\frac12$.
\end{lemma}
\begin{proof}
By definition,
\[
\begin{aligned}S_Nf(x)&=\frac1{2\pi}\int_{-\pi}^\pi f(t)\sum_{|n|\le N}e^{in(x-t)}\dd t\eqbr=\frac1{2\pi}\int_{-\pi}^\pi f(t)D_N(t-x)\dd t,\end{aligned}
\]
using that $D_N$ is even; substitute $t=x+u$ and use periodicity. The closed form is a geometric sum:
$D_N(u)\bigl(e^{iu/2}-e^{-iu/2}\bigr)=e^{i(N+\frac12)u}-e^{-i(N+\frac12)u}$. Finally $\frac1{2\pi}\int_{-\pi}^\pi D_N=1$, since only the term $n=0$
contributes, and the two halves are equal because $D_N$ is even.
\end{proof}

\begin{theorem}[Dirichlet, 1829]\label{thm:dirichletfourier}
Let $f$ be piecewise $C^1$. For every $x$,
\[
S_Nf(x)\xrightarrow[N\to\infty]{}\frac{f(x^+)+f(x^-)}2.
\]
In particular $S_Nf(x)\to f(x)$ wherever $f$ is continuous.
\end{theorem}
\begin{proof}
By Lemma~\ref{lem:dirichletkernel}, writing $\frac12f(x^\pm)=\frac1{2\pi}\int f(x^\pm)D_N$ over $[0,\pi]$ and $[-\pi,0]$ respectively,
\[
\begin{aligned}S_Nf(x)&-\frac{f(x^+)+f(x^-)}2\eqbr=\frac1{2\pi}\int_{-\pi}^\pi G(u)\sin\Bigl(N+\frac12\Bigr)u\dd u,\end{aligned}
\]
where $G(u)=\frac{f(x+u)-f(x^+)}{\sin(u/2)}$ for $0<u\le\pi$ and $G(u)=\frac{f(x+u)-f(x^-)}{\sin(u/2)}$ for $-\pi\le u<0$. Away from $0$, $G$ is
piecewise continuous. Near $0^+$, $f$ coincides on some $\left]x,x+\eta\right]$ with a $C^1$ function $h$ on $[x,x+\eta]$ with $h(x)=f(x^+)$, so
\[
G(u)=\frac{h(x+u)-h(x)}u\cdot\frac u{\sin(u/2)},
\]
which tends to $2h'(x)$ as $u\to0^+$, and similarly at $0^-$. So $G$ is piecewise continuous on $[-\pi,\pi]$, and the integral tends to $0$ by Corollary~\ref{cor:bessel}.
\end{proof}

The proof shows that convergence at $x$ depends only on the behaviour of $f$ near $x$ (the \emph{localisation principle}):
far from $u=0$ the factor $\sin(N+\frac12)u$ oscillates so fast that it averages everything out.

\begin{example}[the sawtooth]\label{ex:sawtooth}
Let $f(x)=\frac{\pi-x}2$ on $\left]0,2\pi\right[$, extended $2\pi$-periodically, with $f(0)=0$. Then $c_0(f)=0$ and, integrating by parts,
$c_n(f)=-\frac i{2n}$ for $n\ne0$, so $S_Nf(x)=\sum_{n=1}^N\frac{\sin nx}n$. The function is piecewise $C^1$, and $\frac{f(0^+)+f(0^-)}2=\frac{\frac\pi2-\frac\pi2}2=0$,
so Theorem~\ref{thm:dirichletfourier} gives, for every real $x$,
\[
\sum_{n\ge1}\frac{\sin nx}n=\frac{\pi-x}2\qquad(0<x<2\pi),
\]
and the sum is $0$ for $x\in2\pi\Z$.
This is the sum promised after Theorem~\ref{thm:dirichlet} in Chapter~1. At $x=\frac\pi2$ it is Leibniz's series
$1-\frac13+\frac15-\cdots=\frac\pi4$ once more.
\end{example}

\section{Fej\'er's theorem}

Pointwise convergence can fail for continuous functions (Exercise~\ref{ex:divergent}). Fej\'er discovered in 1900, at
nineteen, that averaging repairs everything.

\begin{definition}
The \emph{Fej\'er means} of $f$ are $\sigma_Nf=\frac1N\sum_{k=0}^{N-1}S_kf$ for $N\ge1$, trigonometric polynomials of degree at most $N-1$.
\end{definition}

\begin{lemma}[Fej\'er kernel]\label{lem:fejerkernel}
$\sigma_Nf(x)=\frac1{2\pi}\int_{-\pi}^\pi f(x+u)F_N(u)\dd u$, where $F_N=\frac1N\sum_{k<N}D_k$ satisfies, for $u\notin2\pi\Z$,
\[
F_N(u)=\frac1N\Bigl(\frac{\sin(Nu/2)}{\sin(u/2)}\Bigr)^2.
\]
Hence $F_N\ge0$, $\frac1{2\pi}\int_{-\pi}^\pi F_N=1$, and $F_N(u)\le\frac1{N\sin^2(\delta/2)}$ for $\delta\le|u|\le\pi$.
\end{lemma}
\begin{proof}
The integral formula follows from Lemma~\ref{lem:dirichletkernel} by averaging. Since
$2\sin(k+\frac12)u\,\sin\frac u2=\cos ku-\cos(k+1)u$, the sum telescopes:
\[
\sin\frac u2\sum_{k=0}^{N-1}\sin\Bigl(k+\frac12\Bigr)u=\frac{1-\cos Nu}2=\sin^2\frac{Nu}2.
\]
Divide by $N\sin^2\frac u2$. Each $D_k$ has mean $1$, so $F_N$ does too, and $\sin^2\frac u2\ge\sin^2\frac\delta2$ for $\delta\le|u|\le\pi$.
\end{proof}

\begin{figure}[ht]
\centering
\fig{ch06-kernels}
\caption{The Dirichlet kernel $D_8$ (blue) and the Fej\'er kernel $F_8$ (orange) on $[-\pi,\pi]$. Both have mean $1$. The Dirichlet kernel
takes negative values and its oscillations decay slowly; the Fej\'er kernel is non-negative and concentrates at $0$.}
\label{fig:kernels}
\end{figure}

\begin{theorem}[Fej\'er, 1900]\label{thm:fejer}
If $f$ is continuous and $2\pi$-periodic, then $\sigma_Nf\to f$ uniformly on $\R$.
\end{theorem}
\begin{proof}
$f$ is bounded, say by $M$, and uniformly continuous (it is continuous on $[-2\pi,2\pi]$, and periodic): given $\varepsilon>0$ there is
$\delta\in\left]0,\pi\right]$ with $|f(x+u)-f(x)|\le\varepsilon$ for all $x$ and $|u|\le\delta$. Since $F_N\ge0$ has mean $1$, with $\Delta(u)=|f(x+u)-f(x)|$,
\[
\begin{aligned}|\sigma_Nf(x)-f(x)|&\le\frac1{2\pi}\int_{-\pi}^\pi\Delta(u)F_N(u)\dd u\eqbr\le\varepsilon+\frac{2M}{N\sin^2(\delta/2)},\end{aligned}
\]
splitting the integral at $|u|=\delta$. The bound is uniform in $x$, and is at most $2\varepsilon$ for $N$ large.
\end{proof}

The positivity of $F_N$ is the whole difference with $D_N$: an average with non-negative weights of values close to $f(x)$
is close to $f(x)$. Three consequences follow at once.

\begin{corollary}\label{cor:fejer}
\begin{enumerate}[label=(\alph*)]
\item Trigonometric polynomials are uniformly dense among continuous $2\pi$-periodic functions.
\item A continuous $2\pi$-periodic function whose Fourier coefficients all vanish is zero.
\item \emph{Weierstrass's theorem}: every continuous function on $[a,b]$ is a uniform limit of polynomials.
\end{enumerate}
\end{corollary}
\begin{proof}
(a) is Theorem~\ref{thm:fejer}. (b) All $\sigma_Nf$ vanish. (c) By an affine change of variable, take $[a,b]=[-1,1]$. The function
$g(\theta)=f(\cos\theta)$ is continuous, even and $2\pi$-periodic, so $c_{-n}(g)=c_n(g)$ and $\sigma_Ng(\theta)=\sum_na_n\cos n\theta$ for some coefficients $a_n$.
From $\cos(n+1)\theta+\cos(n-1)\theta=2\cos\theta\cos n\theta$, by induction $\cos n\theta=T_n(\cos\theta)$ for a polynomial $T_n$ of degree $n$ (the
Chebyshev polynomials). So $\sup_{[-1,1]}\bigl|f-\sum_na_nT_n\bigr|=\sup_\theta|g-\sigma_Ng|\to0$.
\end{proof}

\section{Parseval's identity}

\begin{theorem}[Parseval]\label{thm:parseval}
For $f$ piecewise continuous, $\|f-S_Nf\|_2\to0$ and
\[
\frac1{2\pi}\int_{-\pi}^\pi|f(t)|^2\dd t=\sum_{n\in\Z}|c_n(f)|^2.
\]
More generally, $\langle f,g\rangle=\sum_nc_n(f)\overline{c_n(g)}$ for $f,g$ piecewise continuous.
\end{theorem}
\begin{proof}
First let $f$ be continuous. $\sigma_{N+1}f$ has degree at most $N$, so by Proposition~\ref{prop:bestapprox}
\[
\begin{aligned}\|f-S_Nf\|_2&\le\|f-\sigma_{N+1}f\|_2\eqbr\le\sup|f-\sigma_{N+1}f|\to0.\end{aligned}
\]
Next let $f$ be piecewise continuous, bounded by $M$, with $k$ points of discontinuity in a period. Given $\eta>0$ small, let $g$
equal $f$ except on intervals of length $\eta$ centred at these points, where $g$ is affine and joins the values of $f$ at the
ends. Then $g$ is continuous, $|f-g|\le2M$, and $\|f-g\|_2^2\le\frac{4M^2k\eta}{2\pi}$, which is small for $\eta$ small. Since $S_Ng$ has
degree at most $N$,
\[
\begin{aligned}\|f-S_Nf\|_2&\le\|f-S_Ng\|_2\eqbr\le\|f-g\|_2+\|g-S_Ng\|_2,\end{aligned}
\]
which is small for $N$ large by the first step. The identity follows from $\|f-S_Nf\|_2^2=\|f\|_2^2-\sum_{|n|\le N}|c_n(f)|^2$. For the
general form, $\langle S_Nf,S_Ng\rangle=\sum_{|n|\le N}c_n(f)\overline{c_n(g)}$ by orthonormality, and by Cauchy--Schwarz
\[
\begin{aligned}\bigl|\langle f,g\rangle-\langle S_Nf,S_Ng\rangle\bigr|&\le\|f-S_Nf\|_2\|g\|_2\eqbr+\|S_Nf\|_2\|g-S_Ng\|_2,\end{aligned}
\]
which tends to $0$.
\end{proof}

Parseval's identity is Pythagoras's theorem with infinitely many orthogonal directions: the $e_n$ form a complete
orthonormal system, and the squared length of $f$ is the sum of the squares of its coordinates. Its physical meaning is
conservation of energy between a signal and its spectrum. Exercise~\ref{ex:x2} uses it for the fourth and fifth proofs of
$\zeta(2)=\frac{\pi^2}6$ and for $\zeta(4)=\frac{\pi^4}{90}$. Tome~7 places it in the theory of Hilbert spaces, where $L^2$ is the completion of
the space of piecewise continuous functions for $\|\cdot\|_2$.

\section{The Gibbs phenomenon}

Let $f$ be the square wave: $f=1$ on $\left]0,\pi\right[$, $f=-1$ on $\left]-\pi,0\right[$, $f(0)=f(\pi)=0$. Then $c_n(f)=\frac2{i\pi n}$ for $n$ odd and $0$ for $n$ even, so
\[
S_{2N-1}f(x)=\frac4\pi\sum_{k=1}^N\frac{\sin(2k-1)x}{2k-1}.
\]
By Theorem~\ref{thm:dirichletfourier} this tends to $f(x)$ for every $x$. In 1898 the physicist Michelson built a machine that summed
$80$ terms of such series and was surprised by bumps near the jumps that did not shrink; Gibbs explained them in 1899
(Wilbraham had found them in 1848).

\begin{theorem}[Gibbs phenomenon]\label{thm:gibbs}
The first maximum of $S_{2N-1}f$ on $\left]0,\pi\right[$ is at $x_N=\frac\pi{2N}$, and
\[
S_{2N-1}f(x_N)\xrightarrow[N\to\infty]{}\frac2\pi\int_0^\pi\frac{\sin t}t\dd t=1.17897\ldots
\]
So the partial sums overshoot the jump of height $2$ by about $0.179$, nearly $9\%$ of the jump, however large $N$ is.
\end{theorem}
\begin{proof}
Differentiating and summing the geometric series of $\cos(2k-1)x$ (the sum $2\sin x\sum_{k\le N}\cos(2k-1)x$ telescopes to $\sin2Nx$),
\[
\begin{aligned}\frac{\mathrm d}{\mathrm dx}S_{2N-1}f(x)&=\frac4\pi\sum_{k=1}^N\cos(2k-1)x\eqbr=\frac2\pi\cdot\frac{\sin2Nx}{\sin x},\end{aligned}
\]
which is positive on $\left]0,\frac\pi{2N}\right[$ and changes sign at $\frac\pi{2N}$: this is the first maximum. With $t_k=\frac{(2k-1)\pi}{2N}$, the midpoints of a
subdivision of $[0,\pi]$ into $N$ intervals of length $\frac\pi N$,
\[
S_{2N-1}f(x_N)=\frac4\pi\sum_{k=1}^N\frac{\sin t_k}{2k-1}=\frac2\pi\cdot\frac\pi N\sum_{k=1}^N\frac{\sin t_k}{t_k},
\]
a Riemann sum of the continuous function $\frac{\sin t}t$ on $[0,\pi]$. It tends to $\frac2\pi\operatorname{Si}(\pi)$, with $\operatorname{Si}(\pi)\approx1.85194$ as in
Figure~\ref{fig:si}.
\end{proof}

\begin{figure}[ht]
\centering
\fig{ch06-gibbs}
\caption{The square wave and its partial sums $S_5$ (green) and $S_{39}$ (blue). The overshoot near the jump moves closer to it but
keeps its height, about $1.179$ (dashed).}
\label{fig:gibbs}
\end{figure}

\begin{think}
Theorem~\ref{thm:dirichletfourier} says $S_Nf(x)\to f(x)$ at every point, and Theorem~\ref{thm:gibbs} says $\sup_{]0,\pi[}S_Nf$ does not tend to $1$. Why is there no
contradiction? What does the Fej\'er mean $\sigma_Nf$ do near the jump?
\end{think}

The phenomenon is universal: a piecewise $C^1$ function with a jump of size $J$ at $a$ is a continuous piecewise $C^1$
function plus $\frac J2$ times a translated square wave, and the first part has uniformly convergent Fourier series
(Exercise~\ref{ex:abs}). So every jump produces an overshoot of $\bigl(\frac1\pi\operatorname{Si}(\pi)-\frac12\bigr)J\approx0.0895J$. Fej\'er means do not overshoot at
all, since $F_N\ge0$ keeps $\sigma_Nf$ between the bounds of $f$; this is why signal processing uses smooth windows instead of
sharp truncations of a spectrum.

\section{Equidistribution}

A real sequence $(u_n)$ is \emph{equidistributed modulo $1$} if, for all $0\le a<b\le1$, the proportion of $n\le N$ with fractional part
$\{u_n\}\in[a,b[$ tends to $b-a$. Fourier analysis on the circle $\R/\Z$ reduces this to exponential sums.

\begin{theorem}[Weyl's criterion, 1916]\label{thm:weyl}
$(u_n)$ is equidistributed modulo $1$ if and only if, for every integer $h\ne0$,
\[
\frac1N\sum_{n=1}^Ne^{2\pi ihu_n}\xrightarrow[N\to\infty]{}0.
\]
\end{theorem}
\begin{proof}
For a $1$-periodic function $F$ let $A_N(F)=\frac1N\sum_{n\le N}F(u_n)$, and say that $F$ is \emph{good} if $A_N(F)\to\int_0^1F$. Good functions form a
vector space. Equidistribution says that the $1$-periodic extensions of the indicators $\mathbf 1_{[a,b[}$ are good, and the
criterion says that the $e^{2\pi ihx}$, $h\ne0$, are good (their integral is $0$); constants are always good.

Suppose that all $\mathbf 1_{[a,b[}$ are good, so all step functions are. A continuous $1$-periodic $F$ is uniformly continuous, so for
$\varepsilon>0$ there is a step function $s$ with $|F-s|\le\varepsilon$ everywhere; then $\bigl|A_N(F)-\int F\bigr|\le2\varepsilon+\bigl|A_N(s)-\int s\bigr|$, so $F$ is good.
In particular the real and imaginary parts of $e^{2\pi ihx}$ are good, which gives the criterion.

Conversely, suppose the criterion holds; then every trigonometric polynomial in $e^{2\pi ihx}$ is good. By Fej\'er's theorem
(applied to $x\mapsto F(\frac x{2\pi})$), every continuous $1$-periodic $F$ is a uniform limit of such polynomials, and the same $2\varepsilon$
argument shows that $F$ is good. Now fix $[a,b[$ and $\varepsilon>0$. There are continuous $1$-periodic functions $F_-\le\mathbf 1_{[a,b[}\le F_+$ with
$\int F_+-\int F_-\le4\varepsilon$: take $F_+$ equal to $1$ on $[a,b]$, $0$ at distance more than $\varepsilon$, affine in between, and $F_-$ similarly
inside $[a,b]$ (or $F_-=0$ if $b-a<2\varepsilon$). Since $A_N$ is monotone,
\[
\begin{aligned}\int F_-&\le\liminf A_N(\mathbf 1_{[a,b[})\eqbr\le\limsup A_N(\mathbf 1_{[a,b[})\le\int F_+,\end{aligned}
\]
and both bounds are within $4\varepsilon$ of $b-a$.
\end{proof}

\begin{corollary}[Bohl, Sierpi\'nski, Weyl, 1909--1910]\label{cor:nalpha}
If $\alpha$ is irrational, $(n\alpha)$ is equidistributed modulo $1$.
\end{corollary}
\begin{proof}
For $h\ne0$, $h\alpha\notin\Z$, so $q=e^{2\pi ih\alpha}\ne1$ and $\bigl|\sum_{n=1}^Nq^n\bigr|=\bigl|\frac{q^{N+1}-q}{q-1}\bigr|\le\frac2{|q-1|}$, a bound independent of $N$.
\end{proof}

\begin{figure}[ht]
\centering
\fig{ch06-weyl}
\caption{The fractional parts $\{n\varphi\}$ for $n=1,\dots,60$, with $\varphi$ the golden ratio, plotted at height $n$. Each new point
falls into one of the largest gaps left by the previous ones; at every stage the gaps take at most three lengths (the
three-gap theorem).}
\label{fig:weyl}
\end{figure}

For rational $\alpha=\frac pq$ the fractional parts take only $q$ values, so irrationality is necessary. Exercise~\ref{ex:benford} applies
the corollary to the first digits of the powers of $2$, and Tome~5, Chapter~11 estimates Weyl sums $\sum e^{2\pi ihf(n)}$ for polynomials $f$ (Weyl differencing, van der Corput), proving the equidistribution of $n^2\alpha$.

\section{Several variables and inversion on \texorpdfstring{$\Z^d$}{Zd}}

For $k\in\Z^d$ and $t\in\R^d$ write $k\cdot t=k_1t_1+\cdots+k_dt_d$. Integrals of continuous functions over the cube $Q=[-\pi,\pi]^d$ are iterated
integrals; by Theorem~\ref{thm:fubinirect} and induction on $d$, the order does not matter. A uniform limit of continuous functions on $Q$
can be integrated term by term, one variable at a time.

\begin{proposition}[Fourier inversion on $\Z^d$]\label{prop:inversionZd}
Let $(a_k)_{k\in\Z^d}$ be complex numbers with $\sum_k|a_k|<\infty$, and $\varphi(t)=\sum_ka_ke^{ik\cdot t}$. The series converges uniformly on $\R^d$,
$\varphi$ is continuous, and for every $k\in\Z^d$
\[
a_k=\frac1{(2\pi)^d}\int_Q\varphi(t)e^{-ik\cdot t}\dd t.
\]
\end{proposition}
\begin{proof}
Enumerate $\Z^d$; the terms are bounded by $|a_k|$, so the series converges uniformly (Weierstrass test, Lemma~\ref{lem:uniform}) and
in any order to the same sum (Lemma~\ref{lem:packets}), and $\varphi$ is continuous. For $k,l\in\Z^d$, the iterated integral factorises:
$\int_Qe^{i(l-k)\cdot t}\dd t=\prod_{j=1}^d\int_{-\pi}^\pi e^{i(l_j-k_j)t_j}\dd t_j$, which is $(2\pi)^d$ if $l=k$ and $0$ otherwise. Multiply the series by $e^{-ik\cdot t}$
(still uniformly convergent) and integrate term by term.
\end{proof}

\begin{example}[return probabilities of random walks]
Let $S_n=X_1+\cdots+X_n$ be a random walk on $\Z^d$ with independent steps distributed like $X$, and let $\varphi(t)=E\bigl[e^{it\cdot X}\bigr]=\sum_kP(X=k)e^{ik\cdot t}$
be its characteristic function. By independence (Tome~3, Chapter~7), $E\bigl[e^{it\cdot S_n}\bigr]=\varphi(t)^n=\sum_kP(S_n=k)e^{ik\cdot t}$, and the
probabilities are summable, so Proposition~\ref{prop:inversionZd} gives
\[
P(S_n=0)=\frac1{(2\pi)^d}\int_Q\varphi(t)^n\dd t.
\]
This is the formula used in Tome~3, Chapter~7, to decide when a walk returns to its start infinitely often.
\end{example}

\section{The isoperimetric inequality}

Among closed curves of given length, which encloses the largest area? The Greeks knew the answer, the circle, but a proof
had to wait for the nineteenth century, and Hurwitz found the shortest one in 1901, with Parseval's identity.

Let $z(s)=x(s)+iy(s)$, $s\in[0,2\pi]$, be a closed curve of length $L$, of class $C^1$, parametrised proportionally to arc length:
$|z'(s)|=\frac L{2\pi}$ for all $s$. If it is simple and traversed counterclockwise, the area it encloses is
\[
\begin{aligned}A&=\frac12\int_0^{2\pi}\bigl(xy'-yx'\bigr)\dd s\eqbr=\frac12\operatorname{Im}\int_0^{2\pi}\overline{z(s)}z'(s)\dd s.\end{aligned}
\]
This is Green's formula; for a polygon it is the shoelace formula, and here we take it as the definition of the area.

\begin{theorem}[isoperimetric inequality]\label{thm:isoperimetric}
$4\pi A\le L^2$, with equality if and only if the curve is a circle.
\end{theorem}
\begin{proof}
Extend $z$ to a $2\pi$-periodic function; it is continuous, and $z'$ is piecewise continuous. Let $c_n=c_n(z)$. Integration by parts,
with $z(-\pi)=z(\pi)$, gives $c_n(z')=inc_n$. By Parseval for $z'$,
\[
\Bigl(\frac L{2\pi}\Bigr)^2=\frac1{2\pi}\int_{-\pi}^\pi|z'|^2=\sum_nn^2|c_n|^2,
\]
and by the general form of Parseval, $\frac1{2\pi}\int z'\bar z=\sum_n(inc_n)\overline{c_n}=i\sum_nn|c_n|^2$, so $A=\pi\sum_nn|c_n|^2$. Hence
\[
L^2-4\pi A=4\pi^2\sum_{n\in\Z}\bigl(n^2-n\bigr)|c_n|^2\ge0,
\]
since $n^2\ge n$ for every integer $n$, with equality only for $n\in\{0,1\}$. So equality holds if and only if $c_n=0$ for $n\ne0,1$,
that is, $z(s)=c_0+c_1e^{is}$: a circle.
\end{proof}

\begin{beauty}
One identity, Pythagoras in a space of functions, gives the energy of a sound, two values of $\zeta$, the shape of a soap
film's boundary and the probability that a random walk comes home. Fourier's ``profound study of nature'' turned out to
be a theorem about orthogonality.
\end{beauty}

\section{Exercises}

\begin{exercise}[$\star$]\label{ex:x2}
Compute the Fourier coefficients of the $2\pi$-periodic functions equal to $x$ and to $x^2$ on $\left]-\pi,\pi\right]$. Deduce $\zeta(2)=\frac{\pi^2}6$ twice,
from Dirichlet's theorem at $x=\pi$ for $x^2$ and from Parseval for $x$, and $\zeta(4)=\frac{\pi^4}{90}$ from Parseval for $x^2$.
\end{exercise}

\begin{exercise}[$\star\star$]\label{ex:abs}
Let $f$ be continuous, $2\pi$-periodic and piecewise $C^1$. Show that $c_n(f')=inc_n(f)$, that $\sum_n|c_n(f)|<\infty$, and that $S_Nf\to f$
uniformly.
\end{exercise}

\begin{exercise}[$\star\star$]\label{ex:heat}
Let $f$ be as in Exercise~\ref{ex:abs}, and for $t>0$ let $u(t,x)=\sum_nc_n(f)e^{-n^2t}e^{inx}$. Show that $\partial_tu=\partial_x^2u$ for $t>0$, and that
$u(t,\cdot)\to f$ uniformly as $t\to0^+$. Which harmonics of the initial temperature survive longest?
\end{exercise}

\begin{exercise}[$\star\star$]\label{ex:wirtinger}
(Wirtinger's inequality.) Let $f$ be real, $2\pi$-periodic and of class $C^1$, with $\int_{-\pi}^\pi f=0$. Show that $\int_{-\pi}^\pi f^2\le\int_{-\pi}^\pi f'^2$, and find
the cases of equality.
\end{exercise}

\begin{exercise}[$\star\star$]\label{ex:benford}
Show that the proportion of $n\le N$ for which $2^n$ starts with the digit $d$ tends to $\log_{10}\bigl(1+\frac1d\bigr)$. How often does $2^n$ start
with $1$? With $7$?
\end{exercise}

\begin{exercise}[$\star\star\star$]\label{ex:divergent}
(A continuous function whose Fourier series diverges at $0$, after Fej\'er.)
\begin{enumerate}[label=(\alph*)]
\item Show that $\bigl|\sum_{j=1}^n\frac{\sin jx}j\bigr|\le1+\pi$ for all $n$ and $x$.
\item Let $P_n(x)=2\sin(2nx)\sum_{j=1}^n\frac{\sin jx}j$. Show that $|P_n|\le2+2\pi$, that $P_n(x)=\sum_{j=1}^n\frac{\cos(2n-j)x-\cos(2n+j)x}j$, and that $P_n(0)=0$.
\item Let $n_k=2^{k^3}$ and $f=\sum_{k\ge1}\frac1{k^2}P_{n_k}$. Show that $f$ is continuous and that $S_{2n_k-1}f(0)\ge k\ln2$.
\end{enumerate}
\end{exercise}

\section{Hints}

\begin{hint}{ex:x2}
Integrate by parts; $c_n(x)=\frac{i(-1)^n}n$ and $c_n(x^2)=\frac{2(-1)^n}{n^2}$ for $n\ne0$.
\end{hint}
\begin{hint}{ex:abs}
Integrate by parts on each piece; the boundary terms cancel by continuity and periodicity. Then Cauchy--Schwarz:
$\sum_{n\ne0}|c_n|=\sum_{n\ne0}\frac1{|n|}\cdot|nc_n|$.
\end{hint}
\begin{hint}{ex:heat}
For $t\ge t_0>0$, the derivatives of the terms are bounded by $n^2|c_n|e^{-n^2t_0}$, which is summable. For the limit, split the series
into $|n|\le M$ and $|n|>M$.
\end{hint}
\begin{hint}{ex:wirtinger}
Parseval for $f$ and $f'$, with $c_0(f)=0$.
\end{hint}
\begin{hint}{ex:benford}
$2^n$ starts with $d$ when $\{n\log_{10}2\}\in\bigl[\log_{10}d,\log_{10}(d+1)\bigr[$. Is $\log_{10}2$ rational?
\end{hint}
\begin{hint}{ex:divergent}
(a) For $0<x<\pi$, treat $j\le\frac1x$ with $|\sin jx|\le jx$, and $j>\frac1x$ with summation by parts (Proposition~\ref{prop:abel}) and
$\bigl|\sum_{j=a}^b\sin jx\bigr|\le\frac1{\sin(x/2)}\le\frac\pi x$. (c) Which frequencies does each $P_{n_k}$ contain?
\end{hint}

\section{Solutions}

\begin{solution}{ex:x2}
\idea Two integrations by parts give the coefficients; then evaluate at the point where the cosines all have sign
$(-1)^n$, or square and integrate.
\writeup For $n\ne0$, $\int_{-\pi}^\pi xe^{-inx}\dd x=\bigl[\frac{xe^{-inx}}{-in}\bigr]_{-\pi}^\pi=\frac{2\pi(-1)^n}{-in}$, so $c_n(x)=\frac{i(-1)^n}n$, and $c_0(x)=0$. Then
$\int x^2e^{-inx}\dd x=\frac2{in}\int xe^{-inx}\dd x$ (the bracket vanishes), so $c_n(x^2)=\frac2{in}\cdot\frac{i(-1)^n}n=\frac{2(-1)^n}{n^2}$, and $c_0(x^2)=\frac{\pi^2}3$. In real form,
\[
x^2=\frac{\pi^2}3+4\sum_{n\ge1}\frac{(-1)^n}{n^2}\cos nx\qquad(|x|\le\pi),
\]
with equality everywhere by Theorem~\ref{thm:dirichletfourier}, since the periodic extension of $x^2$ is continuous and piecewise $C^1$. At $x=\pi$,
$\cos n\pi=(-1)^n$ and $\pi^2=\frac{\pi^2}3+4\zeta(2)$: $\zeta(2)=\frac{\pi^2}6$. Parseval for $x$: $\frac1{2\pi}\int x^2=\frac{\pi^2}3=\sum_{n\ne0}\frac1{n^2}=2\zeta(2)$, the
same value again. Parseval for $x^2$: $\frac1{2\pi}\int x^4=\frac{\pi^4}5=\frac{\pi^4}9+8\zeta(4)$, so $\zeta(4)=\frac{\pi^4}8\bigl(\frac15-\frac19\bigr)=\frac{\pi^4}{90}$.
\end{solution}

\begin{solution}{ex:abs}
\idea Differentiation multiplies the coefficients by $in$; since $f'$ has square-summable coefficients, those of $f$ decay fast
enough to be summable.
\writeup Let $-\pi=a_0<\cdots<a_m=\pi$ be the pieces. On each, $\int_{a_j}^{a_{j+1}}f'e^{-inx}=\bigl[fe^{-inx}\bigr]_{a_j}^{a_{j+1}}+in\int_{a_j}^{a_{j+1}}fe^{-inx}$; since $f$
is continuous the brackets telescope to $f(\pi)e^{-in\pi}-f(-\pi)e^{in\pi}=0$ by periodicity. So $c_n(f')=inc_n(f)$. By Cauchy--Schwarz in $\ell^2$
and Bessel for $f'$,
\[
\begin{aligned}\sum_{n\ne0}|c_n(f)|&\le\Bigl(\sum_{n\ne0}\frac1{n^2}\Bigr)^{1/2}\Bigl(\sum_n|c_n(f')|^2\Bigr)^{1/2}\eqbr\le\sqrt{2\zeta(2)}\,\|f'\|_2.\end{aligned}
\]
By the Weierstrass test $S_Nf$ converges uniformly, to a continuous $g$; integrating term by term, $c_n(g)=c_n(f)$, so $g=f$ by
Corollary~\ref{cor:fejer}(b).
\end{solution}

\begin{solution}{ex:heat}
\idea Each harmonic $e^{inx}$ solves the heat equation once multiplied by $e^{-n^2t}$; the factor $e^{-n^2t}$ also makes every series
converge fast for $t>0$.
\writeup Each term $c_ne^{-n^2t}e^{inx}$ satisfies $\partial_t=\partial_x^2$ (both give $-n^2$ times the term). On $t\ge t_0$, the term and its derivatives in $t$
and in $x$ up to order two are bounded by $(1+n^2)|c_n|e^{-n^2t_0}$, summable, so by Lemma~\ref{lem:uniform}(d) (applied in each variable)
$u$ is differentiable term by term and solves the equation. For the limit, $|u(t,x)-f(x)|\le\sum_n|c_n|(1-e^{-n^2t})$ by Exercise~\ref{ex:abs}.
Given $\varepsilon$, choose $M$ with $\sum_{|n|>M}|c_n|\le\varepsilon$; then the sum is at most $(1-e^{-M^2t})\sum|c_n|+\varepsilon\le2\varepsilon$ for $t$ small. The harmonic
$n$ decays like $e^{-n^2t}$: high frequencies die first, the constant $c_0$, the mean temperature, survives, and after a while
only $c_0+c_{\pm1}e^{-t}e^{\pm ix}$ is visible.
\end{solution}

\begin{solution}{ex:wirtinger}
\idea On the Fourier side, differentiation multiplies the $n$-th coefficient by $in$, and $|n|\ge1$ once $c_0$ is gone.
\writeup By Exercise~\ref{ex:abs}, $c_n(f')=inc_n(f)$, and $c_0(f)=0$. By Parseval,
$\frac1{2\pi}\int f'^2=\sum n^2|c_n|^2\ge\sum|c_n|^2=\frac1{2\pi}\int f^2$. Equality holds iff $c_n=0$ for $|n|\ge2$, that is, $f(x)=a\cos x+b\sin x$.
\end{solution}

\begin{solution}{ex:benford}
\idea The first digit of $2^n$ is read from the fractional part of $\log_{10}2^n=n\log_{10}2$, and these fractional parts are equidistributed.
\writeup $2^n$ starts with $d$ iff $d\cdot10^m\le2^n<(d+1)10^m$ for some $m$, iff $\{n\log_{10}2\}\in[\log_{10}d,\log_{10}(d+1)[$. If $\log_{10}2=\frac pq$ with $q\ge1$,
then $2^q=10^p$, impossible since $5$ does not divide $2^q$ (and $p=0$ gives $2^q=1$). So by Corollary~\ref{cor:nalpha} the proportion tends
to $\log_{10}(d+1)-\log_{10}d=\log_{10}(1+\frac1d)$: about $30.1\%$ of powers of $2$ start with $1$, and $5.8\%$ with $7$. This is Benford's law,
observed in many tables of data that span several orders of magnitude.
\end{solution}

\begin{solution}{ex:divergent}
\idea Build blocks of trigonometric polynomials that stay bounded but whose partial sums at $0$ are as large as $\ln n$, then
place the blocks at frequencies so far apart that they do not interact.
\writeup (a) The sum is odd and $2\pi$-periodic in $x$, so take $0<x<\pi$, and let $m=\lfloor\frac1x\rfloor$. For $j\le m$, $\bigl|\frac{\sin jx}j\bigr|\le x$, so these
terms contribute at most $mx\le1$. For $j>m$, the partial sums $\sum_{j=m+1}^b\sin jx=\frac{\cos(m+\frac12)x-\cos(b+\frac12)x}{2\sin(x/2)}$ are at most $\frac1{\sin(x/2)}\le\frac\pi x$
in absolute value (since $\sin\frac x2\ge\frac x\pi$ on $[0,\pi]$). By Abel summation with the decreasing weights $\frac1j$, as in the proof of
Theorem~\ref{thm:dirichlet}, the rest contributes at most $\frac1{m+1}\cdot\frac\pi x\le\pi$, since $m+1>\frac1x$.

(b) The bound follows from (a). The product formula $2\sin a\sin b=\cos(a-b)-\cos(a+b)$ gives the expansion, whose frequencies lie in
$[n,3n]$; and $P_n(0)=0$ since $\sin0=0$.

(c) The series converges uniformly, by (b) and $\sum\frac1{k^2}<\infty$, so $f$ is continuous and its coefficients are obtained term by
term. Since $3n_k<n_{k+1}$, the blocks occupy disjoint ranges of frequencies. The partial sum $S_{2n_k-1}$ keeps all the blocks
$i<k$ (their frequencies are at most $3n_i<n_k$), none of the blocks $i>k$ (frequencies at least $n_{k+1}>2n_k$), and from block $k$ exactly the terms
$\frac{\cos(2n_k-j)x}{j}$. At $x=0$, the blocks $i<k$ give $\frac{P_{n_i}(0)}{i^2}=0$, so
\[
S_{2n_k-1}f(0)=\frac1{k^2}\sum_{j=1}^{n_k}\frac1j\ge\frac{\ln n_k}{k^2}=k\ln2\to\infty.
\]
So the Fourier series of the continuous function $f$ diverges at $0$. In 1966 Carleson proved that this can only happen on a set
of measure zero; Kahane and Katznelson showed that any set of measure zero can occur.
\end{solution}

\section*{Sources and further reading}

\begin{itemize}
\item T.~W.~K\"orner, \emph{Fourier Analysis}, Cambridge University Press, 1988. Short chapters, full of history and
applications; the Gibbs phenomenon, Weyl and the heat equation are all there.
\item E.~Stein and R.~Shakarchi, \emph{Fourier Analysis: An Introduction}, Princeton University Press, 2003, Chapters~2--4
(including Hurwitz's isoperimetric proof and Weyl's criterion).
\item Y.~Katznelson, \emph{An Introduction to Harmonic Analysis}, 3rd ed., Cambridge University Press, 2004.
\item A.~Hurwitz, \emph{Sur quelques applications g\'eom\'etriques des s\'eries de Fourier}, Annales scientifiques de l'\'E.N.S.
19 (1902), 357--408.
\item H.~Weyl, \emph{\"Uber die Gleichverteilung von Zahlen mod.~Eins}, Mathematische Annalen 77 (1916), 313--352.
\item E.~Hewitt and R.~Hewitt, \emph{The Gibbs--Wilbraham phenomenon: an episode in Fourier analysis}, Archive for History of
Exact Sciences 21 (1979), 129--160.
\end{itemize}
